\subsection[Planck's Self-Dual Identity on the Periodicity of Order 108]{Planck self-dual identity in the circular realization of the temporal endomorphism of order 108}
\label{sec:identidad-autodual-planck-ct108-rev3}
\index{Planck!self-dual identity}
\index{calendar of order 108! circular realization}
\index{Compton!reduced and complete readings}

The action line constructed in the preceding section admits a
circular realization compatible with temporal periodicity of order 108.
This realization coordinates three objects that already possess their own types: the
action section, the geometry of one revolution, and the gravitational
transduction. The resulting coordination provides an internal Planck
identity and prepares the positive line on which the general mass
operator will subsequently act.

\subsubsection{Oriented action letters and declared circular realization}

Let
\[
 \mathcal L_S^+,\quad \mathcal L_T^+,\quad \mathcal L_C^+,
 \quad \mathcal L_L^+,\quad \mathcal L_M^+,\quad \mathcal L_G^+
\]
the positive half-lines of action, time, velocity, length, and mass.
The pre-return and post-return coordinates are written
\begin{equation}
 \hbar_{\rm pre}=H_5\,10^{-34}\mathcal U_S,
 \qquad
 \hbar_{\rm ret}=\eta_{\rm ret}\,10^{-34}\mathcal U_S,
 \qquad
 R_{\rm act}=\frac{\hbar_{\rm pre}}{\hbar_{\rm ret}}
 =\frac{H_5}{\eta_{\rm ret}}.
 \label{eq:planck-ct108-cartas-pre-ret-rev3}
\end{equation}
The subscript \(a\in\{\mathrm{pre},\mathrm{ret}\}\) shall designate either
of the two charts. In each one the exact distinction is preserved
\begin{equation}
 h_a=2\pi\hbar_a .
 \label{eq:planck-ct108-h-hbar-rev3}
\end{equation}

Let us fix an elementary duration \(t_0\in\mathcal L_T^+\), an internal
velocity \(c_{\rm int}\in\mathcal L_C^+\), and
\(\ell_0=c_{\rm int}t_0\). The declared circular realization of the
temporal periodicity of order 108 is determined by
\begin{equation}
 T_{108}=108t_0,
 \qquad
 \omega_{108}=\frac{2\pi}{T_{108}},
 \qquad
 R_{108}=\frac{c_{\rm int}}{\omega_{108}}
 =\frac{54}{\pi}\ell_0,
 \label{eq:planck-ct108-realizacion-circular-rev3}
\end{equation}
and by the length of the full turn
\begin{equation}
 C_{108}=2\pi R_{108}=c_{\rm int}T_{108}=108\ell_0.
 \label{eq:planck-ct108-perimetro-rev3}
\end{equation}
Thus, \(R_{108}\) is the radius of the realization and \(C_{108}\) its
perimeter. The angular frequency, the cyclic frequency and the two
action constants are related through
\(\omega_{108}=2\pi/T_{108}\) and \(h_a=2\pi\hbar_a\).

\subsubsection{Correspondence between the order-108 periodicity and Compton readings}

In the chart \(a\), the circular quantum and its mass section are
\begin{equation}
 E_{108,a}=\hbar_a\omega_{108},
 \qquad
 m_{108,a}=\frac{E_{108,a}}{c_{\rm int}^{2}}
 =\frac{\hbar_a\omega_{108}}{c_{\rm int}^{2}}.
 \label{eq:planck-ct108-cuanto-circular-rev3}
\end{equation}

\begin{theorem}[Identity between circular periodicity and Compton readings]
\label{thm:planck-ct108-compton-rev3}
For each chart \(a\in\{\mathrm{pre},\mathrm{ret}\}\), the reduced and complete
Compton readings satisfy
\begin{equation}
 \boxed{
 \bar\lambda_{\rm C}(m_{108,a})
 =\frac{\hbar_a}{m_{108,a}c_{\rm int}}=R_{108},
 \qquad
 \lambda_{\rm C}(m_{108,a})
 =\frac{h_a}{m_{108,a}c_{\rm int}}=C_{108}.}
 \label{eq:planck-ct108-identidad-compton-rev3}
\end{equation}
\end{theorem}

\begin{proof}
Substituting
\eqref{eq:planck-ct108-cuanto-circular-rev3} produces
\[
 \frac{\hbar_a}{m_{108,a}c_{\rm int}}
 =\frac{c_{\rm int}}{\omega_{108}}=R_{108}.
\]
The second equality follows by applying simultaneously
\eqref{eq:planck-ct108-h-hbar-rev3} and
\eqref{eq:planck-ct108-perimetro-rev3}.
\end{proof}

The theorem is homogeneous in
\((\hbar_a,c_{\rm int},t_0)\): the same periodicity determines the radius,
the perimeter, and the two Compton charts. The circular realization is the
declared geometric primitive; the preceding equalities are its
exact consequences.

\subsubsection{Circular selector and internal Planck unit}

The gravitational line is coordinated with the chart \(a\) through the homogeneous
transducer
\begin{equation}
 G_a(\theta)=\theta\,
 \frac{c_{\rm int}^{5}t_0^2}{\hbar_a}
 \in\mathcal L_G^+,
 \qquad \theta\in\mathbb R_{>0}.
 \label{eq:planck-ct108-transductor-rev3}
\end{equation}
Its reduced gravitational radius, evaluated on the circular quantum, is
\begin{equation}
 r_{{\rm g},a}(\theta)
 :=\frac{G_a(\theta)m_{108,a}}{c_{\rm int}^{2}}
 =\theta\frac{\pi}{54}\ell_0.
 \label{eq:planck-ct108-radio-gravitatorio-rev3}
\end{equation}

\begin{proposition}[Planck Circular Selector]
\label{prop:selector-circular-planck-ct108-rev3}
Within the circular realization
\eqref{eq:planck-ct108-realizacion-circular-rev3}, the condition
\(r_{{\rm g},a}(\theta)=R_{108}\) selects a unique positive coordinate,
\begin{equation}
 \boxed{\theta_{108}^{\rm circ}=\left(\frac{54}{\pi}\right)^2.}
 \label{eq:theta-circular-planck-ct108-rev3}
\end{equation}
The coordinate is independent of the chart oriented \(a\).
\end{proposition}

\begin{proof}
The equations
\eqref{eq:planck-ct108-realizacion-circular-rev3} and
\eqref{eq:planck-ct108-radio-gravitatorio-rev3} reduce the condition to
\[
 \theta\frac{\pi}{54}\ell_0=\frac{54}{\pi}\ell_0.
\]
The positivity of \(\ell_0\) and the strict monotonicity of the map
\(\theta\mapsto r_{{\rm g},a}(\theta)\) provide the solution and its
uniqueness.
\end{proof}

Let \(G_{108,a}^{\rm circ}:=G_a(\theta_{108}^{\rm circ})\). The internal units
associated with this chart satisfy
\begin{align}
 \ell_{{\rm P},a}^{\rm circ}
 &:=\sqrt{\frac{\hbar_a G_{108,a}^{\circ}}{c_{\rm int}^{3}}}
 =R_{108},
 &
 t_{{\rm P},a}^{\circ}
 &:=\sqrt{\frac{\hbar_a G_{108,a}^{\circ}}{c_{\rm int}^{5}}}
 =\omega_{108}^{-1},
 \label{eq:planck-ct108-longitud-tiempo-rev3}\\
 m_{{\rm P},a}^{\circ}
 &:=\sqrt{\frac{\hbar_a c_{\rm int}}{G_{108,a}^{\circ}}}
 =m_{108,a},
 &
 E_{{\rm P},a}^{\circ}
 &:=m_{{\rm P},a}^{\circ}c_{\rm int}^{2}=E_{108,a}.
 \label{eq:planck-ct108-masa-energia-rev3}
\end{align}
The equality of radii is therefore published as
\begin{equation}
 R_{108}=\bar\lambda_{\rm C}=r_{\rm g}=\ell_{\rm P}^{\rm circ},
 \qquad
 C_{108}=\lambda_{\rm C}=2\pi r_{\rm g}=2\pi\ell_{\rm P}^{\rm circ}.
 \label{eq:planck-ct108-tres-lectores-rev3}
\end{equation}

The convention of this theorem is the reduced gravitational radius
\(r_{\rm g}=Gm/c_{\rm int}^{2}\). The Schwarzschild radius satisfies
\(r_{\rm s}=2r_{\rm g}\); if this second convention is adopted, its crossing with
\(\bar\lambda_{\rm C}\) occurs at
\(m=m_{{\rm P},a}^{\circ}/\sqrt2\). The declaration of the radius therefore forms
part of the chart and prevents a factor of two from being inadvertently transferred
between the two identities.

The coordinate \(\theta_{108}^{\rm circ}\) belongs to the declared circular
specialization. The general gravitational selector, formulated on the
enriched holonomy and its spectral response, retains its residence in
\cref{sec:selector-gravitatorio-rev11}. Thus, the circular
identity fixes an internal Planck normalization and the general criterion
fixes the conditions that the physical realization of the gravitational
section must satisfy.

\subsubsection{Pre-return/post-return covariance}

The product of the reciprocal sections of action and gravity preserves the
circular geometry:
\begin{equation}
 \frac{m_{108,\rm pre}}{m_{108,\rm ret}}=R_{\rm act},
 \qquad
 \frac{G_{108,\rm pre}^{\circ}}{G_{108,\rm ret}^{\circ}}
 =R_{\rm act}^{-1},
 \qquad
 \hbar_a G_{108,a}^{\circ}
 =\theta_{108}^{\rm circ}c_{\rm int}^{5}t_0^2.
 \label{eq:planck-ct108-covariancia-pre-ret-rev3}
\end{equation}
In particular, \(\ell_{{\rm P},\rm pre}^{\rm circ}\) and
\(\ell_{{\rm P},\rm ret}^{\rm circ}\) coincide with \(R_{108}\). The
bidirectional memory resides in the quotient of the masses and in its
gravitational reciprocal; the geometric scale resides in the invariant product.

Reciprocity admits a symmetric parametrization that separates the common
scale from the orientation of the return. Define
\begin{equation}
 \begin{aligned}
 \hbar_{\rm geom}
 &:=\sqrt{\hbar_{\rm pre}\hbar_{\rm ret}},
 &
 G_{108,\rm geom}^{\circ}
 &:=\sqrt{G_{108,\rm pre}^{\circ}G_{108,\rm ret}^{\circ}},
 &
 \xi_{\rm act}&:=\frac12\log R_{\rm act}.
 \end{aligned}
 \label{eq:planck-ct108-parametrizacion-simetrica-rev3}
\end{equation}
Then
\[
 \hbar_{\rm pre}=\hbar_{\rm geom}e^{\xi_{\rm act}},
 \qquad
 \hbar_{\rm ret}=\hbar_{\rm geom}e^{-\xi_{\rm act}},
\]
whereas the conjugate gravitational sections satisfy
\[
 G_{108,\rm pre}^{\circ}
 =G_{108,\rm geom}^{\circ}e^{-\xi_{\rm act}},
 \qquad
 G_{108,\rm ret}^{\circ}
 =G_{108,\rm geom}^{\circ}e^{\xi_{\rm act}}.
\]
The coordinate \(\xi_{\rm act}\) preserves the pre-return/post-return
orientation; the product \(\hbar_aG_{108,a}^{\circ}\)
preserves geometric scale. This parametrization of the action charts
does not by itself introduce distinct propagation speeds.

\subsubsection{Split self-dual mass line}

Given a chart \(a\), let the split algebra be
\[
 \mathbb D_{\rm sp}:=\mathbb R[j]/(j^2-1),
 \qquad
 P_\pm:=\frac{I\pm j}{2},
\]
and consider its spectral decomposition
\[
 \mathbb D_{\rm sp}
 =\mathbb R P_+\oplus\mathbb R P_-,
 \qquad
 P_\pm^2=P_\pm,\quad
 P_+P_-=0,\quad
 P_++P_-=I.
\]
The broken conjugation swaps the projectors,
\[
 \overline{uP_++vP_-}=vP_++uP_-,
\]
and its norm is
\[
 N_{\rm sp}(uP_++vP_-):=uv.
\]

For \(m\in\mathcal L_M^+\), define the section
\begin{equation}
 Z_a(m)
 :=\frac{m_{{\rm P},a}^{\circ}}{m}P_+
   +\frac{m}{m_{{\rm P},a}^{\circ}}P_- .
 \label{eq:planck-ct108-seccion-partida-masa-rev3}
\end{equation}
The unit \(m_{{\rm P},a}^{\circ}\) also determines the involution
\begin{equation}
 \iota_{{\rm P},a}(m)
 =\frac{(m_{{\rm P},a}^{\circ})^2}{m}.
 \label{eq:planck-ct108-involucion-masas-rev3}
\end{equation}

\begin{theorem}[Split self-dual mass line]
\label{thm:planck-ct108-recta-partida-masa-rev3}
For every \(m\in\mathcal L_M^+\),
\begin{equation}
 \boxed{
 N_{\rm sp}(Z_a(m))=1,
 \qquad
 \overline{Z_a(m)}
 =Z_a\!\left(\iota_{{\rm P},a}(m)\right).}
 \label{eq:planck-ct108-norma-conjugacion-rev3}
\end{equation}
Conjugation has a unique positive fixed point,
\[
 m=m_{{\rm P},a}^{\circ},
 \qquad
 Z_a(m_{{\rm P},a}^{\circ})=I.
\]
\end{theorem}

\begin{proof}
The norm is the product of the two reciprocal coordinates:
\[
 \frac{m_{{\rm P},a}^{\circ}}m
 \frac m{m_{{\rm P},a}^{\circ}}=1.
\]
Conjugation interchanges both coordinates, which is equivalent to replacing
\(m\) by \((m_{{\rm P},a}^{\circ})^2/m\). The positive fixed point satisfies
\(m^2=(m_{{\rm P},a}^{\circ})^2\).
\end{proof}

With the speed
\begin{equation}
 x_a(m)=\log\frac{m}{m_{{\rm P},a}^{\circ}},
 \label{eq:planck-ct108-rapidez-masa-rev3}
\end{equation}
the section and its conjugate are written
\[
 Z_a(m)=e^{-x_a}P_++e^{x_a}P_-,
 \qquad
 \overline{Z_a(m)}
 =e^{x_a}P_++e^{-x_a}P_-.
\]
Their even and odd readers are
\begin{equation}
 Q_a(m):=\frac{e^{x_a}+e^{-x_a}}2=\cosh x_a,
 \qquad
 A_a(m):=\frac{e^{x_a}-e^{-x_a}}2=\sinh x_a.
 \label{eq:planck-ct108-lectores-par-impar-rev3}
\end{equation}
The first preserves the distance to the self-dual point and the second preserves its orientation.

The reciprocal lattice vectors associated with the same coordinate are
\[
 r_{{\rm g},a}(m)=\frac{G_{108,a}^{\circ}m}{c_{\rm int}^{2}},
 \qquad
 \bar\lambda_{{\rm C},a}(m)=\frac{\hbar_a}{mc_{\rm int}},
\]
and satisfy
\begin{equation}
 \frac{r_{{\rm g},a}(m)}{\bar\lambda_{{\rm C},a}(m)}
 =e^{2x_a(m)},
 \qquad
 \frac{\bar\lambda_{{\rm C},a}(m)}{r_{{\rm g},a}(m)}
 =e^{-2x_a(m)}.
 \label{eq:planck-ct108-lectores-exponenciales-rev3}
\end{equation}
Conjugation interchanges the two readers; Planck's identity is their unitary equilibrium.

\begin{exactbox}[title={Planck identity and electron closure}]
The identity \(m=m_{{\rm P},a}^{\circ}\) fixes the self-dual origin of the positive line of masses. The electron section fixes a distinct position on that same line through its own genealogy. The \cref{sec:ley-general-masa-accion-autonoma} constructs the universal operator that transports absolute scale, character, towers, and bidirectional reading; the \cref{sec:cierre-electron-two-way-rev11} then realizes its central rank-one reduction. Thus the Planck unit provides the identity of the coordinate, and the electron provides the first complete material transition on it.
\end{exactbox}

\begin{proposition}[Genealogical realization of the electronic center]
\label{prop:planck-ct108-realizacion-centro-electron-rev3}
Let
\[
 F_e^+=\mathbb R_{>0}(P_++P_-)
\]
the positive half-line of the rank-one central fiber, let
\(\mathcal G_e\) be the space of enriched electronic
genealogies and let \(\mathcal L_M^+\) be the positive mass
line. A realization morphism
\[
 \mathfrak R_e:
 F_e^+\times\mathcal G_e\longrightarrow\mathcal L_M^+
\]
can send the identity \(I=P_++P_-\) to a mass different from
\(m_{{\rm P},a}^{\circ}\) without altering the
\cref{thm:planck-ct108-recta-partida-masa-rev3}: the Planck identity fixes the
autodual origin of the chart, while the electronic genealogy fixes a
displacement on that chart.
\end{proposition}

\begin{proof}
If \(\chi_e:\mathcal G_e\to\mathbb R_{>0}\) is the positive character constructed
by the electronic record, the rule
\[
 \mathfrak R_e(\lambda I,g)
 =\lambda\,m_{{\rm P},a}^{\circ}\chi_e(g)
\]
is positive and multiplicative in the genealogical coordinate. For the neutral
genealogy, \(\chi_e=1\), Planck's origin is recovered; for a genealogy with
\(\chi_e(g)\ne1\), the same internal identity occupies a different position on
the line. The specific electronic closure preserves its source in the
\cref{sec:cierre-electron-two-way-rev11}.
\end{proof}

\paragraph{Domain of the Planck self-dual identity: temporal periodicity of order \(108\), declared circular realization, and separation of the electronic closure}
The pre-return/post-return action line and the temporal periodicity of
order (108) precede the circular realization, which is taken as a declared
geometric primitive. The identity between circular
periodicity and the Compton readings, the
selector \(\theta_{108}^{\rm circ}\), the internal Planck unit, and the
involution of the positive line are exact results once that
realization has been fixed. General gravitational selection and electronic
closure are constructed separately in the indicated sections.
